\section{Experiments}
\label{sec:experiments}

To evaluate whether phased orchestration improves task completion quality, we design four experiments that collectively address the concerns raised in prior review: (1) an ablation study isolating each phase's contribution, (2) a control marker effectiveness study, (3) a three-protocol consistency evaluation, and (4) a cost-benefit analysis. All experiments are conducted through the Pigs HTTP API, requiring only API access---no GPU resources.
为评估分阶段编排是否能提升任务完成质量，我们设计了四组实验，以共同回应先前评审中提出的关注点：（1）隔离各阶段贡献的消融研究，（2）控制标记有效性研究，（3）三协议一致性评估，（4）成本效益分析。所有实验均通过 Pigs HTTP API 执行，仅需 API 访问——无需 GPU 资源。

\subsection{Experimental Setup}

\textbf{System configuration.} All experiments use the Pigs API server with default orchestration limits (\texttt{max\_post\_iterations=3}, \texttt{max\_pre\_replans=2}), continuation TTL of 30 minutes, capacity of 256, and Chinese (zh) phase prompts. Thinking effort is set to xhigh for OpenAI and max for Anthropic.
\textbf{系统配置。} 所有实验使用 Pigs API 服务器，采用默认编排限制（\texttt{max\_post\_iterations=3}、\texttt{max\_pre\_replans=2}），续传 TTL 为 30 分钟，容量为 256，阶段提示采用中文（zh）。思考力度对 OpenAI 设为 xhigh，对 Anthropic 设为 max。

\textbf{Benchmarks.} We use three public benchmarks:
\textbf{基准测试。} 我们使用三个公开基准测试：

\begin{center}
\begin{tabular}{lllr}
\toprule
Benchmark & Domain & Tasks & Public \\
\midrule
SWE-bench Lite~\citep{jimenez2024swebench} & Software engineering & 300 & \checkmark \\
GAIA~\citep{mialon2023gaia} & General AI assistant & 300 & \checkmark \\
AgentBoard~\citep{ma2024agentboard} & Multi-turn agent & 9 tasks & \checkmark \\
\bottomrule
\end{tabular}
\end{center}

\textbf{Ablation conditions.} For Experiment 1, we define four conditions:
\textbf{消融条件。} 对于实验一，我们定义四种条件：
\begin{itemize}[nosep]
    \item \textbf{Condition A (Baseline)}: Direct API call without orchestration (non-\texttt{-pig} model, passthrough).
    \textbf{条件 A（基线）}：不经过编排的直接 API 调用（非 \texttt{-pig} 模型，透传）。
    \item \textbf{Condition B (Full)}: Complete Pre$\rightarrow$Executor$\rightarrow$Post orchestration (\texttt{-pig} model).
    \textbf{条件 B（完整）}：完整的 Pre$\rightarrow$Executor$\rightarrow$Post 编排（\texttt{-pig} 模型）。
    \item \textbf{Condition C (Pre-only)}: Only the Pre phase (planning without execution verification).
    \textbf{条件 C（仅 Pre）}：仅 Pre 阶段（规划但不执行验证）。
    \item \textbf{Condition D (Executor+Post)}: Skip Pre, directly execute and review.
    \textbf{条件 D（Executor+Post）}：跳过 Pre，直接执行并审查。
\end{itemize}

\textbf{Metrics.} We measure: (1) task completion rate (benchmark-specific), (2) average LLM API calls per task, (3) total input + output tokens per task, (4) wall-clock latency, and (5) error recovery rate (percentage of tasks where PIGFAIL triggered a successful replan).
\textbf{指标。} 我们测量：（1）任务完成率（基准测试特定），（2）每任务平均 LLM API 调用次数，（3）每任务总输入 + 输出令牌数，（4）端到端延迟，（5）错误恢复率（\texttt{PIGFAIL} 触发成功重新规划的任务百分比）。

\subsection{Experiment 1: Ablation Study}

\textbf{Setup.} We evaluate four conditions (A: baseline, B: full orchestration, C: Pre-only, D: Executor+Post) on SWE-bench Lite (300 tasks) and GAIA (300 tasks). Each task is sent as a \texttt{-pig} model request (Condition B/C/D) or a passthrough request (Condition A) to the Pigs API.
\textbf{设置。} 我们在 SWE-bench Lite（300 个任务）和 GAIA（300 个任务）上评估四种条件（A：基线，B：完整编排，C：仅 Pre，D：Executor+Post）。每个任务作为 \texttt{-pig} 模型请求（条件 B/C/D）或透传请求（条件 A）发送至 Pigs API。

\textbf{Procedure.} For each task:
\textbf{流程。} 对于每个任务：
\begin{enumerate}[nosep]
    \item Send the task prompt to the Pigs API endpoint.
    将任务提示发送至 Pigs API 端点。
    \item If tool calls are returned, execute them locally and send results back.
    若返回工具调用，在本地执行并将结果发回。
    \item Repeat until the API returns a final response (no tool calls).
    重复直至 API 返回最终响应（无工具调用）。
    \item Evaluate the final response against the benchmark's ground truth.
    将最终响应与基准测试的标准答案进行评估。
\end{enumerate}

\textbf{Hypothesis.} Full orchestration (B) outperforms baseline (A) on task completion rate, with the Pre phase contributing planning and the Post phase contributing verification.
\textbf{假设。} 完整编排（B）在任务完成率上优于基线（A），其中 Pre 阶段贡献规划能力，Post 阶段贡献验证能力。

\subsection{Experiment 2: Control Marker Effectiveness}

\textbf{Setup.} We compare two configurations on GAIA: (A) with \texttt{PIGFAIL} enabled (failure triggers replan), (B) with \texttt{PIGFAIL} disabled (failure treated as completion).
\textbf{设置。} 我们在 GAIA 上比较两种配置：（A）启用 \texttt{PIGFAIL}（失败触发重新规划），（B）禁用 \texttt{PIGFAIL}（失败被视为完成）。

\textbf{Procedure.} We count: (1) tasks where PIGFAIL was emitted, (2) tasks where replan led to successful completion, (3) tasks that failed despite replan.
\textbf{流程。} 我们统计：（1）发出 \texttt{PIGFAIL} 的任务数，（2）重新规划后成功完成的任务数，（3）尽管重新规划仍然失败的任务数。

\textbf{Hypothesis.} PIGFAIL-based replan improves outcomes on tasks where the initial approach was flawed.
\textbf{假设。} 基于 \texttt{PIGFAIL} 的重新规划在初始方法有缺陷的任务上能改善结果。

\subsection{Experiment 3: Three-Protocol Consistency}

\textbf{Setup.} We send the same set of 100 tasks through three API protocols: OpenAI Chat (\texttt{/chat/completions}), Anthropic Messages (\texttt{/v1/messages}), and OpenAI Responses (\texttt{/responses}).
\textbf{设置.} 我们将同一组 100 个任务通过三种 API 协议发送：OpenAI Chat（\texttt{/chat/completions}）、Anthropic Messages（\texttt{/v1/messages}）和 OpenAI Responses（\texttt{/responses}）。

\textbf{Procedure.} For each task and protocol, we record: (1) whether the task completed, (2) the visible text output, (3) any tool calls made. We measure consistency as the percentage of tasks where all three protocols agree on completion status.
\textbf{流程.} 对于每个任务和协议，我们记录：（1）任务是否完成，（2）可见文本输出，（3）所执行的工具调用。我们将一致性定义为三种协议在完成状态上达成一致的任务百分比。

\textbf{Hypothesis.} The protocol-native request preservation ensures consistent behavior across all three protocols.
\textbf{假设.} 协议原生请求保留确保了三种协议间行为的一致性。

\subsection{Experiment 4: Cost-Benefit Analysis}

\textbf{Setup.} On SWE-bench Lite, we measure per-task: (1) total input tokens, (2) total output tokens, (3) number of API calls, (4) wall-clock latency, (5) task completion.
\textbf{设置.} 在 SWE-bench Lite 上，我们逐任务测量：（1）总输入令牌数，（2）总输出令牌数，（3）API 调用次数，（4）端到端延迟，（5）任务完成情况。

\textbf{Procedure.} We compare Condition A (baseline) vs.\ Condition B (full orchestration) and compute the overhead ratio for each metric.
\textbf{流程.} 我们比较条件 A（基线）与条件 B（完整编排），并计算各指标的开销比率。

\textbf{Hypothesis.} The token overhead of orchestration (3 phases vs.\ 1) is justified by improved completion rate.
\textbf{假设.} 编排的令牌开销（3 阶段 vs.\ 1 阶段）可通过提升的完成率来证明其合理性。
